\section{Final checklist}

This chapter introduced discrete random variables as functions on sample spaces.
The central workflow was
\[
\Omega\longrightarrow X(\omega)\longrightarrow P(X=x).
\]
The sample space describes possible outcomes of the experiment. The random
variable extracts a numerical quantity. The distribution describes how
probability is transferred from outcomes to values.

\subsection*{Checklist for constructing a discrete random variable}

When a problem involves a discrete random variable, proceed as follows.

\begin{enumerate}
    \item Identify the original experiment.
    \item Construct or recall the sample space \(\Omega\).
    \item Define the random variable as a function
    \[
    X:\Omega\to\mathbb{R}.
    \]
    \item State clearly what \(X(\omega)\) means for an elementary outcome
    \(\omega\).
    \item Find the support:
    \[
    \operatorname{supp}(X)=\{x:P(X=x)>0\}.
    \]
    \item For each support value \(x\), identify the event
    \[
    \{X=x\}=\{\omega\in\Omega:X(\omega)=x\}.
    \]
    \item Compute the probability mass function
    \[
    p_X(x)=P(X=x).
    \]
    \item Check that
    \[
    \sum_x p_X(x)=1.
    \]
\end{enumerate}

\subsection*{Checklist for interpreting the distribution}

After the probability law has been found, ask:

\begin{enumerate}
    \item What does each value of \(X\) mean in the original experiment?
    \item Which values are most likely?
    \item Is the support finite or countably infinite?
    \item Does the mass function come from counting, weighting, or both?
    \item What is the cumulative distribution function
    \[
    F_X(t)=P(X\leq t)?
    \]
    \item How do jumps of the CDF correspond to masses of the PMF?
\end{enumerate}

\subsection*{Checklist for numerical summaries}

For expectation and variance, use the distribution of \(X\):
\[
\E[X]=\sum_x xP(X=x),
\]
\[
\Var(X)=\E[(X-\E[X])^2].
\]
For computation, one may also use
\[
\Var(X)=\E[X^2]-(\E[X])^2.
\]

Interpret these quantities carefully:
\begin{itemize}
    \item expectation is a weighted average or center of mass;
    \item expectation need not be a possible value of \(X\);
    \item variance measures average squared deviation from the expectation;
    \item standard deviation is the square root of variance and uses the original
    units of the variable.
\end{itemize}

\subsection*{Checklist for classical laws}

Before applying a named distribution, identify the mechanism.

\begin{itemize}
    \item Bernoulli: one success-or-failure trial.
    \item Binomial: number of successes in a fixed number of independent trials.
    \item Geometric: waiting time until the first success.
    \item Hypergeometric: number of successes in sampling without replacement.
    \item Poisson: count of events in a fixed observation window.
\end{itemize}

Do not select a distribution because a formula looks familiar. Select it because
the assumptions of the mechanism match the experiment.

\begin{summarybox}
A discrete random variable is a numerical view of an experiment. The probability
law of that variable is obtained by collecting all elementary outcomes that
produce the same value and adding their probabilities. Classical distributions
are reusable models of common mechanisms, not isolated formulas.
\end{summarybox}
